\section{Elementary Number Theory}
\label{sec:elementary-number-theory}

\subsection{Primes and Divisibility}
\label{sec:primes-and-divisibility}

\begin{definition}[Divisors]
    Given \(a, b \in \ZZ\), we say \emph{\(a\) divides \(b\)} if there exists \(c \in \ZZ\) such that \(b = ac\).

    Equivalently, we might also say \emph{\(a\) is a divisor of \(b\)}, \emph{\(a\) is a factor of \(b\)}, or \emph{\(b\) is a multiple of \(a\)}. We denote \(a \divides b\).
\end{definition}

\begin{definition}[Proper Factors]
    For any \(b \in \ZZ\), \(\para{\pm 1}\) and \(\para{\pm b}\) are always factors of \(b\). We call those the \emph{trivial factors} of \(b\), and the rest \emph{proper/non-trivial factors}.
\end{definition}

\begin{definition}[Prime, Composite]
    A natural number \(n \geq 2\) is \emph{prime} if its only divisors are trivial. Otherwise, we say \(n\) is \emph{composite}.
\end{definition}

\begin{proposition}[Fundamental Theorem of Arithmetics, Part 1]
    Every natural number \(\geq 2\) can be written as a product of primes.
\end{proposition}

\begin{proof}
    We show this by induction on \(n\). This is true for \(n = 2\).

    Now, let \(n > 2\) and suppose the claim holds up to and including \(n - 1\). If \(n\) is prime, then we are done. If \(n\) is composite, then by definition, \(n = ab\) for some \(1 < a, b < n\).

    By the induction hypothesis, \(a = p_1 p_2 \cdots p_k\) and \(b = q_1 q_2 \cdots q_l\) for some primes \(p_1, p_2, \ldots, p_k\) and \(q_1, q_2, \ldots, q_l\).

    Hence,
    \[
        n = ab = p_1 p_2 \cdots p_k q_1 q_2 \cdots q_l
    \]
    is a product of primes.
\end{proof}

\begin{theorem}[Infinitely Many Primes]
    There are infinitely many primes.
\end{theorem}

\begin{proof}[(Euclid, 300 BC)]
    Suppose for the sake of contradiction, that there are only finitely many primes, and let them be \(p_1, p_2, \ldots, p_k\).

    Consider \(N = p_1 p_2 \cdots p_k + 1\). \(N\) is not itself a prime. Then \(p_1 \notdivides N\), or otherwise,
    \[
        p_1 \divides \para{N - p_1 p_2 \cdots p_k} = 1.
    \]

    Likewise, \(p_2 \notdivides N, \ldots, p_k \notdivides N\).

    This contradicts with the fact that \(N\) can be written as a product of primes.
\end{proof}

The previous proposition states the existence of a representation as a product of primes, but our proof does not give uniqueness. It is in fact unique, and the following lemma is useful.

\begin{proposition}[Euclid's Lemma]
    If \(p\) is prime and \(p \divides ab\), then \(p \divides a\) and \(p \divides b\).
\end{proposition}

Before we prove this, we introduce what is meant by the highest common factor.

\begin{definition}[Highest Common Factor/Greatest Common Divisor]
    Given \(a, b \in \NN\), we say a number \(c \in \NN\) is the \emph{highest common factor} (hcf) or \emph{greatest common divisor} (gcd) of \(a\) and \(b\) if:
    \begin{enumerate}
        \item \(c \divides a\) and \(c \divides b\) (i.e. \(c\) is a common divisor of \(a\) and \(b\)); and
        \item \(d \divides a\) and \(d \divides b\) implies \(d \divides c\) (i.e. every common divisor of \(a\) and \(b\) divides \(c\)).
    \end{enumerate}

    We denote this as
    \[
        c = \hcf\para{a, b} = \gcd\para{a, b} = \para{a, b}.
    \]
\end{definition}

\begin{example}[Greatest Common Divisor]
    We attempt to find the greatest common divisor of \(12\) and \(18\).

    The factors of \(12\) are \(1, 2, 3, 4, 6, 12\).

    The factors of \(18\) are \(1, 2, 3, 6, 9, 18\).

    So the common factors of \(12\) and \(18\) are \(1, 2, 3, 6\).

    Hence,
    \[
        \gcd\para{12, 18} = 6.
    \]
\end{example}

This definition does not immediately lead to the fact that the greatest common divisor always exists -- if it does it must be the 'greatest' common factor, but it is not trivial that the 'greatest' common factor is always a multiple of every other common factor.

We show by construction that in fact it must always exist, but we first introduce an important method.

\begin{proposition}[Division Algorithm]
    Let \(n, k \in \NN\), then we can write \(n = qk + r\) for some \(q, r \in \ZZ\), with \(0 \leq r \leq k - 1\).
    
    Such \(q\) is called the \emph{quotient} and \(r\) is called the \emph{remainder}.
\end{proposition}

\begin{proof}
    We show this by induction on \(n\). This is true for \(n = 1\).

    Suppose \(n \geq 2\), and the statement holds for \(n - 1\), i.e.
    \[
        n - 1 = qk + r
    \]
    for some \(q, r \in \ZZ\) and \(0 \leq r \leq k - 1\).

    If \(0 \leq r < k - 1\), then
    \[
        n = qk + \para{r + 1}
    \]
    and hence \(1 \leq r + 1 < k\) is within the required range.

    If \(r = k - 1\), then
    \[
        n = qk + \para{k - 1} + 1 = \para{q + 1}k + 0,
    \]
    and the obtained \(q\) and \(r\) are within the range.
\end{proof}

\begin{remark}
    The \(q\) and \(r\) obtained from this integer division is unique. Indeed, if \(n = qk + r = q'k + r'\), then
    \[
        \para{q - q'}k = r' - r
    \]
    and we see that since \(0 \leq r, r' < k\), we have
    \[
        -k < r' - r < k
    \]
    and hence since \(\para{q - q'}k\) is a multiple of \(k\), we must have \(r - r' = 0\), and \(q - q' = 0\).
\end{remark}

We now use this integer division to conduct an algorithm called \emph{Euclid's algorithm}, as in \autoref{tab:euclids-algorithm}.

\begin{table}
    \centering
    \begin{tabular}{c|cc|c}
        Input & \(a, b\) && \(a = 372, b = 162\) \\
        \hline
        Step 1 & \(a = q_1 b + r_1\) & \(q_1, r_1 \in \ZZ, 0 \leq r_1 \leq b - 1\) & \(372 = 2 \cdot \boxed{162} + \boxed{48}\) \\
        Step 2 & \(b = q_2 r_1 + r_2\) & \(q_2, r_2 \in \ZZ, 0 \leq r_2 \leq r_1 - 1\) & \(162 = 3 \cdot \boxed{48} + \boxed{18}\) \\
        Step 3 & \(r_1 = q_3 r_2 + r_3\) & \(q_3, r_3 \in \ZZ, 0 \leq r_3 \leq r_2 - 1\) & \(48 = 2 \cdot \boxed{18} + \boxed{2}\) \\
        \(\vdots\) & \(\vdots\) & \(\vdots\) & \\
        Step \(n\) & \(r_{n - 2} = q_n r_{n - 1} + r_n\) & \(q_n, r_n \in \ZZ, 0 \leq r_n \leq r_{n - 1} - 1\) & \(18 = 1 \times \boxed{12} + \boxed{6}\) \\
        Step \(n + 1\) & \(r_{n - 1} = q_{n + 1} r_n + r_{n + 1}\) & \(q_{n + 1} \in \ZZ, \boxed{r_{n + 1} = 0}\) & \(12 = 2 \times \boxed{6} + \boxed{0}\) \\
        \hline
        Output & \(r_n\) & & \(6\)
    \end{tabular}
    \caption{Euclid's algorithm}
    \label{tab:euclids-algorithm}
\end{table}

Note that this algorithm terminates in \(n < b\) steps, since
\[
    b > r_! > r_2 > \cdots > r_n > 0.
\]

\begin{theorem}[Euclid's Algorithm gives the Greatest Common Divisor]
    The output of Euclid's algorithm with input \(a, b\) is \(\gcd\para{a, b}\).
\end{theorem}

\begin{proof}
    We show that it satisfies the two defining properties of the greatest common divisor.

    \begin{enumerate}
        \item We must have \(r_n \divides r_{n - 1}\), as \(r_{n + 1} = 0\) at step \(n + 1\). Therefore, \(r_N \divides r_{n - 2}\) by step \(n\).
        
        Hence, inductively, \(r_n \divides r_i\) for \(i = 1, 2, \ldots, n - 1\), and hence \(r_n \divides b\) by step 2, and \(r_n \divides a\) by step 1.

        \item Given any \(d\) such that \(d \divides a\) and \(d \divides b\). From step 1 we see \(d \divides r_1\), and from step 2 we have \(d \divides r_2\).
        
        Hence, inductively, \(d \divides r_i\) for all \(i = 1, 2, \ldots, n\), and in particular, \(d \divides r_n\).
    \end{enumerate}
\end{proof}

\begin{example}[Greatest Common Divisor of \(87\) and \(52\)]
    Suppose we want to find \(\gcd\para{87, 52}\). We run Euclid's Algorithm:
    \begin{align*}
        87 &= 1 \cdot 52 + 35, \\
        52 &= 1 \cdot 35 + 17, \\
        35 &= 2 \cdot 17 + \boxed{1}, \\
        17 &= 17 \cdot \boxed{1},
    \end{align*}
    so \(\gcd\para{87, 52} = 1\). Note we may terminate the algorithm one step early, outputting the final non-zero divisor, which is the greatest common divisor.
\end{example}

\begin{definition}[Coprime]
    When \(\gcd\para{a, b} = 1\), we say \(a\) and \(b\) are \emph{coprime}.
\end{definition}

Notice that we may reverse Euclid's algorithm. For example, in the previous example,
\begin{align*}
    1 &= 35 - 2 \times 17 \\
       &= 35 - 2 \times \para{52 - 1 \times 35} \\
       &= 3 \times 35 - 2 \times 52 \\
       &= 3 \times \para{57 - 1 \times 52} - 2 \times 52 \\
       &= 3 \times 57 - 5 \times 52.
\end{align*}

This leads to the following lemma.

\begin{theorem}[B\'{e}zout's Identity]
    For all \(a, b \in \NN\), there exists \(x, y \in \ZZ\), such that
    \[
        ax + by = \gcd\para{a, b}.
    \]

    We can write \(\gcd\para{a, b}\) as a linear combination of \(a\) and \(b\).
\end{theorem}

\begin{proof}[by reversing Euclid's Algorithm]
    We run Euclid's Algorithm with input \(a, b\) to obtain an output \(r_n\). We have
    \[
        r_n = x r_{n - 1} + y r_{n - 2}
    \]
    for some \(x, y \in \ZZ\).

    For step \(n - 1\), we see that \(r_{n - 1}\) is expressible of some linear combination of \(r_{n - 2}\) and \(r_{n - 3}\). We have
    \[
        r_{n - 1} = x r_{n - 2} + y r_{n - 3}
    \]
    for some \(x, y \in \ZZ\).

    Hence,
    \[
        r_n = x r_{n - 2} + y r_{n - 3}
    \]  
    for some \(x, y \in \ZZ\).

    We may continue inductively, for all \(i = 1, 2, \ldots, n - 1\), to see that
    \[
        r_n = x r_i + y r_{i - 1}
    \]
    for some \(x, y \in \ZZ\). Hence,
    \[
        r_{n} = xa + yb
    \]
    for some \(x, y \in \ZZ\) from steps 1 and 2.
\end{proof}

\begin{remark}
    Euclid's Algorithm not only proves the existence of such \(x, y \in \ZZ\), but also gives a quick way to find them.
\end{remark}

\begin{proof}[by utilising well--ordering principle]
    Let \(h\) be the least positive linear combination of \(a\) and \(b\), i.e. the least positive integer of the form \(ax + by\) for some \(x, y \in \ZZ\). Such \(h\) exists by well--ordering principle.

    We aim to show that \(h = \gcd\para{a, b}\).

    To show the first property, suppose that \(h\) does not divide one of \(a\) or \(b\), say \(a\). Then by integer division, we have
    \[
        a = qh + r
    \]
    for some \(q, r \in \ZZ\) with \(0 < r < h\).

    Hence,
    \begin{align*}
        r &= a - qh \\
        &= a - q \cdot \para{xa + yb} \\
        &= \para{1 - qx} a - \para{qy} b
    \end{align*}
    is a positive linear combination of \(a\) and \(b\), strictly smaller than \(h\), and contradicting with the fact that \(h\) is minimal.

    To show the second property, notice that for any common divisor \(d \divides a\) and \(d \divides b\), we have \(d \divides ax + by\) for all \(x, y \in \ZZ\), and \(d \divides h\) in particular.
\end{proof}

\begin{remark}
    This proof tells us that \(\gcd\para{a, b}\) exists and is a linear combination of \(a\) and \(b\), but it gives in fact no way to find the coefficients \(x, y \in \ZZ\) to find \(\gcd\para{a, b}\).
\end{remark}

A natural question to ask, is by changing the right-hand side, when does an equation of the form \(ax + by = c\) have a solution? For example, is there a solution in \(\para{x, y} \in \ZZ \times \ZZ\) to
\[
    160 x + 72 y = 33?
\]

What about
\[
    87 x + 52 y = 33?
\]

It is obvious for the second equation, that we have for \(x', y' \in \ZZ\) that
\[
    87 x' + 52 y' = 1
\]
so let \(x = 33 x'\) and \(y = 33 y'\), we have
\[
    87 x + 52 y = 33.
\]

\begin{corollary}[B\'{e}zout's Identity, Continued]
    Let \(a, b \in \NN\). Then the equation
    \[
        ax + by = c
    \]
    has a solution for \(x, y \in \ZZ\) if and only if
    \[
        \gcd\para{a, b} \divides c.
    \]
\end{corollary}

\begin{proof}
    Let \(h = \gcd\para{a, b}\)

    \(\implies\). Suppose there are such integers \(x, y \in \ZZ\) such that \(ax + by = c\). Since \(h \divides a\) and \(h \divides b\), we must have \(h \divides ax + by = c\).

    \(\impliedby\). Conversely, suppose \(h \divides c\). Then by the previous theorem, there exists \(x, y \in \ZZ\), such that \(h = ax + by\).

    But then \(c = \frac{ch} \cdot h\) and \(\frac{c}{h} \in \ZZ\), and hence
    \[
        c = \frac{c}{h} \cdot \para{ax + by} = a \cdot \para{x \cdot \frac{c}{h}} + b \cdot \para{y \cdot \frac{c}{h}}
    \]
    constructing a solution as desired.
\end{proof}

We now prove Euclid's Lemma.

\begin{proposition}[Euclid's Lemma]
    If \(p\) is a prime and \(p \divides ab\), then \(p \divides a\) or \(p \divides b\).
\end{proposition}

\begin{proof}
    Suppose that \(p \divides ab\) but \(p \notdivides a\). We want to prove that \(p \divides b\).

    Since \(p\) is prime and \(p \notdivides a\), we have \(\gcd\para{a, p} = 1\).

    Then, by B\'{e}zout's Identity, there exists \(x, y \in \ZZ\) such that
    \[
        xp + ya = 1.
    \]

    Hence, multiplying both sides by \(b\) gives
    \[
        \para{xb} \cdot p + y \cdot \para{ab} = b,
    \]
    and we have \(p \divides \para{xb} \cdot b\) and \(p \divides y \cdot \para{ab}\). Hence, \(p \divides b\).
\end{proof}

\begin{remarks}
    \begin{enumerate}
        \item We can iterate this lemma: if \(p \divides a_1 a_2 \cdots a_n\), then \(p \divides a_i\) for some \(i = 1, 2, \ldots, n\).
        \item We do require \(p\) is prime. Indeed, \(4 \divides 2 \times 6\), but \(4 \notdivides 2\) and \(4 \notdivides 6\).
    \end{enumerate}
\end{remarks}

\begin{theorem}[Fundamental Theorem of Arithemetic]
    Every natural number \(n \geq 2\) is expressible as a product of primes uniquely, up to reordering. 
\end{theorem}

\begin{proof}
    The existence of such proposition follows from the earlier proposition.

    To show uniqueness, we show by strong induction. This is indeed true for \(n = 2\).

    Consider \(n > 2\). Suppose that there are two factorisations
    \begin{align*}
        n &= p_1 p_2 \cdots p_k,\\
        n &= q_1 q_2 \cdots q_l
    \end{align*}
    where \(p_i\) and \(q_j\) are all prime. We want to prove \(k = l\), and after reordering \(p_i = q_i\) for all \(i = 1, 2, \ldots, k\).

    Since
    \[
        p_1 p_2 \cdots p_k = q_1 q_2 \cdots q_l,
    \]
    we have
    \[
        p_1 \divides n = q_1 q_2 \cdots q_l
    \]
    and by Euclid's Lemma, we must have
    \[
        p_1 \divides q_i
    \]
    for some \(i\), and after relabelling, we assume \(p_1 \divides q_1\).

    Since \(q_1\) is prime, it follows that \(p_1 = q_1\). Hence,
    \[
        \frac{n}{p_1} = p_2 \cdots p_k = q_2 \cdots q_l < n
    \]
    and by the strong induction hypothesis, in fact \(k = l\), and after reordering, \(p_2 = q_2, p_3 = q_3, \ldots, p_k = q_k\).
\end{proof}

\begin{remark}
    There are 'arithmetical systems' (rings) permitting addition and multiplication in which factorisation is not unique. For example, consider the ring extension
    \[
        \ZZ\brac{\sqrt{-3}} = \set{x + y i \sqrt{3}}{x, y \in \ZZ}
    \]

    We can add and multiply to get another element of \(\ZZ\brac{\sqrt{-3}}\), e.g.
    \begin{align*}
        \para{1 + \sqrt{-3}} + \para{1 - \sqrt{-3}} &= 2,\\
        \para{1 + \sqrt{-3}} \para{1 - \sqrt{-3}} &= 1^2 - \para{-3} = 4.
    \end{align*}

    In \(\ZZ\brac{\sqrt{-3}}\) we may define what it means to be \emph{prime}, and both \(\para{1 \pm \sqrt{3}}\) are indeed prime.

    However, \(4 = 2 \cdot 2\) and \(2\) is a prime in this system as well.

    Factorisation is \emph{not} unique in this system.
\end{remark}

We may now apply the fundamental theorem of arithmetics to some useful applications.

\begin{definition}[Least Common Multiple]
    Given \(a, b \in \NN\), we say a number \(c \in \NN\) is the \emph{least common multiple} (lcm) of \(a\) and \(b\) if:
    \begin{enumerate}
        \item \(a \divides c\) and \(b \divides c\) (i.e. \(c\) is a common multiple of \(a\) and \(b\)); and
        \item \(a \divides m\) and \(b \divides m\) implies \(c \divides m\) (i.e. every common multiple of \(a\) and \(b\) is a multiple of \(c\)).
    \end{enumerate}

    We denote this as
    \[
        c = \lcm\para{a, b} = \brac{a, b}.
    \]  
\end{definition}

\begin{examples}[Applications of Fundamental Theorem of Arithmetics]
    \begin{enumerate}
        \item What are the factors of \(n = 2^3 \cdot 3^7 \cdot 11\)? All number of the form \(2^a \cdot 3^b \cdot 11^c\) for \(0 \leq a \leq 3, 0 \leq b \leq 7\) and \(0 \leq c \leq 1\), and there must not be any others.
        
        For example, if \(7\) divides a factor of \(n\), then \(7\) divides \(n\), and hence we have a factorisation of \(n\) involving \(7\), contradicting with the uniqueness of the factorisation.

        More generally, the factors of
        \[
            n = p_1^{\alpha_1} p_2^{\alpha_2} \cdots p_k^{\alpha_k}
        \]
        are \emph{precisely} the numbers of the form
        \[
            p_1^{\beta_1} p_2^{\beta_2} \cdots p_2^{\beta_k}
        \]
        where \(0 \leq \beta_i \leq \alpha_i\) for all \(i = 1, 2, \ldots, k\).

        \item What are the common factors of \(2^3 \cdot 3^7 \cdot 5 \cdot 11^2\) and \(2^4 \cdot 3^2 \cdot 11 \cdot 13\)? All numbers of the form \(2^a \cdot 3^b \cdot 11^c\) for \(0 \leq a \leq 3, 0 \leq b \leq 2\) and \(0 \leq c \leq 1\). Then the greatest common divisor is \(2^3 \cdot 3^2 \cdot 11\).
        
        In general, the greatest common divisor of two numbers
        \[
            p_1^{\alpha_1} p_2^{\alpha_2} \cdots p_k^{\alpha_k}
        \]
        and
        \[
            p_1^{\beta_1} p_2^{\beta_2} \cdots p_k^{\beta_k}
        \]
        for \(\alpha_i, \beta_i \geq 0\), is
        \[
            p_1^{\min\brce{\alpha_1, \beta_1}} p_2^{\min\brce{\alpha_2, \beta_2}} \cdots p_k^{\min\brce{\alpha_k, \beta_k}}.
        \]

        \item What are the common multiples of the two previous numbers? They are all numbers of the form
        \[
            2^a \cdot 3^b \cdot 5^c \cdot 11^d \cdot 13^e
        \]
        for \(a \geq 4, b \geq 7, c \geq 1, d \geq 3, e \geq 1\), times any integer.

        Hence, the number \(2^4 \cdot 3^7 \cdot 5 \cdot 11^3 \cdot 13\) is a common multiple, and every other common multiple is a multiple of it, and hence it is the least common multiple of the two numbers.

        In general, the least common multiple of two numbers
        \[
            p_1^{\alpha_1} p_2^{\alpha_2} \cdots p_k^{\alpha_k}
        \]
        and
        \[
            p_1^{\beta_1} p_2^{\beta_2} \cdots p_k^{\beta_k}
        \]
        for \(\alpha_i, \beta_i \geq 0\), is
        \[
            p_1^{\max\brce{\alpha_1, \beta_1}} p_2^{\max\brce{\alpha_2, \beta_2}} \cdots p_k^{\max\brce{\alpha_k, \beta_k}}.
        \]
    \end{enumerate}
\end{examples}

\begin{proposition}[GCD and LCM multiple to the Product]
    We have for any \(x, y \in \NN\), that
    \[
        \gcd\para{x, y} \cdot \lcm\para{x, y} = x \cdot y.
    \]
\end{proposition}

\begin{proof}
    This is immediate from the previous example and using the fact that
    \[
        \min\brce{a, b} + \max\brce{a, b} = a + b.
    \]
\end{proof}

We may now present another proof that there are infinitely many primes --- recall that we did not use the fact that there are infinitely many primes in any part of the form.

\begin{proof}[of infinitely may primes, Erd\"{o}s, 1930]
    Let \(p_1, p_2, \ldots, p_k\) be all the primes. Since all numbers is uniquely expressed as a product of primes, consider
    \[
        p_1^{j_1} p_2^{j_2} \cdots p_k^{j_k}.
    \]

    We may rewrite this in the following way:
    \[
        p_1^{j_1} p_2^{j_2} \cdots p_k^{j_k} = m^2 p_1^{i_1} p_2^{i_2} \cdots p_k^{i_k}
    \]
    for \(i_l = 0, 1\).

    Now, let \(N \in \NN\). Suppose given a number \(\leq N\) of this form, we have \(m \leq \sqrt{N}\), and hence we have at most \(\sqrt{N} \cdot 2^k\) numbers of this form that are \(\leq N\).

    Now, take \(N > \sqrt{N} \cdot 2^k\), i.e. \(N > 4^k\), then there must be a number which is \(\leq N\) but not of this form, not having a prime factorisation, which gives a contradiction.
\end{proof}

\begin{remark}
    Both of these proofs in fact give an upper bound on the \(k\)th prime. Euclid's Proof tells us that the \(k\)th prime, \(p_k\) has an upper bound
    \[
        p_k < 2^{2^k}.
    \]

    Even better, Erd\"{o}s's Proof gives the bound
    \[
        p_k < 4^k.
    \]

    In fact, a much tighter bound is given by the \emph{prime number theorem}, stating
    \[
        p_k \sim k \log k.
    \]
\end{remark}

\subsection{Modular Arithmetic}
\label{sec:modular-arithmetic}

\begin{definition}[Integers Modulo \(n\)]
    Let \(n \geq 2\) be a natural number. Then, the \emph{integers modulo \(n\)}, denoted \(\ZZ_{n}\) or \(\ZZ / n\ZZ\), consists of the integers with the two integers being the same if they differ by a multiple of \(n\).

    Formally, if \(x\) and \(y\) are the same in \(\ZZ_n\), we write
    \[
        x \equiv y \pmod{n}
    \]
    or
    \[
        x \equiv y \quad \para{n}
    \]
    or
    \[
        x = y \text{ in } \ZZ_n.
    \]

    This is equivalent to \(n \divides x - y\), or \(x = y + kn\) for some \(k \in \ZZ\).
\end{definition}

\begin{example}[Example of Congurence]
    In \(\ZZ_7\), \(2\) is regarded the same as \(16\), and hence
    \[
        2 \equiv 16 \pmod 7.
    \]
\end{example}

Instead of a straight axis, we can view \(\ZZ_n\) as a circle, as in \autoref{fig:visuallisation-of-zn}.

\begin{figure}[ht!]
    \centering
    \caption{Visualisation of \(\ZZ_N\)}
    \label{fig:visuallisation-of-zn}
\end{figure}

\begin{proposition}[Addition and Multiplication on \(\ZZ_n\)]
    Addition and multiplication on \(\ZZ_n\) is well defined.
\end{proposition}

\begin{proof}
    If \(a \equiv a', b \equiv b' \pmod{n}\), then
    \[
        n \divides \para{a - a'} + \para{b - b'} = \para{a + b} - \para{a' + b'}
    \]
    and hence
    \[
        a + b \equiv a' + b' \pmod{n}.
    \]

    Similarly,
    \[
        n \divides \para{a - a'} b + a' \cdot \para{b - b'} = ab - a' b'
    \]
    and hence
    \[
        ab \equiv a' b' \pmod{n}.
    \]
\end{proof}

This allows us to do usual arithmetic modulo \(n\).

We first see a useful application of modular arithmetic, that is determining whether a Diophantine equation has a solution.

\begin{examples}[Diophantine Equation]
    Does \(2a^2 + 3b^3 = 1\) have a solution in \(\ZZ\)? If there is one, then we have
    \[
        2a^2 \equiv 1 \pmod{3}.
    \]

    However, note that
    \begin{align*}
        2 \cdot 0^2 &\equiv 0, \\
        2 \cdot 1^2 &\equiv 2, \pmod{3} \\
        2 \cdot 2^2 &\equiv 2,
    \end{align*}
    and \(1 \nequiv 2 \pmod{3}\).

    Therefore, there is no solution.
\end{examples}

